% RENDU-3D.tex
% Du terrain nu au rendu HR 3D : comment une carte se construit et se dessine.
% Compilation, en double passe :
%   cd docs/cartes && lualatex RENDU-3D.tex && lualatex RENDU-3D.tex
\documentclass[a4paper,12pt]{article}
\usepackage[left=1.5cm,right=1.5cm,top=2cm,bottom=2.5cm]{geometry}
\usepackage[utf8]{inputenc}
\usepackage[T1]{fontenc}
\usepackage[french]{babel}

\usepackage{iftex}
\ifPDFTeX
  \usepackage{fourier}
\else
  \usepackage{fourier-otf}
\fi

\usepackage{amsmath}
% Pas d'amssymb : fourier fournit déjà les symboles AMS (\mathbb compris) et
% les deux entrent en conflit.
\usepackage{hyperref}
\hypersetup{colorlinks=true, linkcolor=blue, urlcolor=blue}

% Noms de fichiers et de fonctions : \code{...} plutôt que \texttt partout,
% pour pouvoir changer la présentation en un seul endroit.
\newcommand{\code}[1]{\texttt{#1}}

\setlength{\parindent}{0pt}
\setlength{\parskip}{0.5\baselineskip}

% Les longs noms de code ne se coupent pas : on laisse les espaces s'étirer
% pour renvoyer le nom entier à la ligne plutôt que de déborder dans la marge.
\sloppy
\setlength{\emergencystretch}{3em}

% Version affichée sous le titre : celle de la notice (docs/version.tex,
% régénéré par build.sh -d à partir du dernier tag). En son absence
% (compilation manuelle), valeur de repli.
\InputIfFileExists{../version.tex}{}{}
\providecommand{\noticeversion}{version de développement}

\begin{document}

\begin{center}
  {\LARGE\bfseries Du terrain nu au rendu HR 3D}\\[0.5ex]
  {\large Comment une carte d'Artouste se construit et se dessine}\\[1ex]
  Olivier Booklage\\
  version \noticeversion{}, le \today
\end{center}

Quatre étapes : partir de zéro, le rendu LR, le rendu HR, le rendu HR 3D. Les
fonctions sont citées à chaque étape pour retrouver le code.

\section{Fabriquer une carte}

Une carte se fabrique hors du jeu, avec
\code{python3 tools/fetch\_terrain.py <zone>}. Ce script en appelle quatre
autres :

\begin{itemize}
  \item \code{tools/terrain/zones/<zone>.py} décrit la zone : emprise en
    lon/lat, taille de la grille d'altitudes, hauteur de l'orthophoto, point
    de départ, hélisurfaces ;
  \item \code{tools/terrain/relief.py} demande les altitudes à l'API
    altimétrie de l'IGN (RGE ALTI,
    \url{https://data.geopf.fr/altimetrie/1.0/calcul/alti/rest/elevation.json})
    et les écrit dans \code{heightmap.png}, un PNG 16 bits étalé de 0 à
    \code{elev\_max} ;
  \item \code{tools/terrain/ortho.py} télécharge l'orthophoto BD ORTHO par le
    WMS de la Géoplateforme (\url{https://data.geopf.fr/wms-r/wms}, couche
    \code{ORTHOIMAGERY.ORTHOPHOTOS}), comble le blanc en bordure de couverture
    (\code{fill\_nodata()}), donne une couleur unie à la mer
    (\code{flatten\_sea()}) et écrit \code{ortho.jpg} ;
  \item \code{tools/terrain/outputs.py} et \code{meta.py} écrivent
    \code{terrain.txt} (emprise, maillage, altitude maximale). Le script pose
    aussi \code{landmarks.txt} (lieux remarquables) et \code{helipads.txt}.
\end{itemize}

Le dossier \code{assets/terrain/<carte>/} obtenu est la carte livrée : l'état
LR. Toutes les données viennent de l'IGN
(\url{https://geoservices.ign.fr/}), sous Licence Ouverte Etalab 2.0.

\section{Le rendu LR : un maillage et une photo}

Au chargement, \code{Terrain::Terrain()} (\code{src/render/Terrain.cpp}) lit
\code{terrain.txt}, \code{heightmap.png} et \code{ortho.jpg}. Puis
\code{Terrain::construireMaillage()}
(\code{src/render/terrain/TerrainMaillage.cpp}) échantillonne la heightmap
sous un plafond de sommets (\code{relief\_sommets\_max}) : sur une carte de
18~km, la maille tombe à 15-20~m. Le maillage est construit une fois,
\code{Terrain::draw()} le dessine à chaque image.

Le shader (\code{assets/shaders/terrain.frag}) drape \code{ortho.jpg} dessus
et y mêle un grain rocheux selon la pente. La physique lit le même relief par
\code{Terrain::heightAt()} (interpolation bilinéaire de la heightmap).

Résultat : correct de loin, flou au ras du sol (une dizaine de mètres par
pixel), anguleux sur les crêtes.

\newpage

\section{Le rendu HR : des tuiles d'image fines}

Le gestionnaire de cartes (touche \code{C}, puis Entrée) fabrique des tuiles
d'orthophoto à 0,25~m par pixel. \code{Fabrique::lancer()}
(\code{src/app/cartes/FabriqueTuiles.cpp}) télécharge l'image bloc par bloc
(même WMS, même couche que le socle), la découpe en tuiles de 512~px, les
compresse en BC7 et les écrit en \code{.dds} (\code{render::dds::ecrire()}),
avec un index par niveau de finesse. Le script
\code{tools/terrain/fetch\_tuiles.py} fait pareil en ligne de commande ; les
deux produisent le même jeu et une fabrication commencée par l'un se reprend
par l'autre.

Au chargement, \code{ouvrirNiveaux()} (\code{src/render/tuiles/Pyramide.cpp})
retrouve le jeu. Une \code{Fenetre} (\code{src/render/tuiles/Fenetre.hpp}) par
niveau garde une texture de tuiles autour de la caméra :
\code{Fenetre::suivre()} la recentre à chaque image et charge les tuiles qui
entrent, un masque dit au shader lesquelles sont prêtes. Dans
\code{terrain.frag}, \code{poserTuiles()} pose le niveau large puis le niveau
serré sur l'orthophoto d'ensemble, avec un fondu par distance
(\code{reglerNiveau()} dans
\code{src/app/render/ApplicationRenderTerrain.cpp}).

Résultat : le sol est net en vol bas. Le relief, lui, reste celui du maillage
LR.

\section{Le rendu HR 3D : la fenêtre de relief}

La méthode est propre à Artouste. Elle n'est pas reprise d'un moteur existant,
elle a été construite pour ce jeu, mesure par mesure. Ses quatre choix : le
relief fin est un \emph{écart} ajouté à la carte, pas une surface de plus
(carte + détail = laser, par construction) ; le pas des tuiles s'emboîte dans
la maille du terrain ; la finesse ne dépend que de la distance, pour que les
grilles lisent la même surface à leur jonction ; la physique recalcule la
formule du shader, donc l'appareil se pose sur ce que le pilote voit.

La touche \code{L} du gestionnaire fabrique le jeu de relief.
\code{Fabrique::lancerRelief()} puis \code{Fabrique::boucleRelief()}
(\code{src/app/cartes/fabrique/FabriqueReliefBoucle.cpp}) téléchargent le MNT
LiDAR HD de l'IGN, le modèle du sol nu (couche
\code{IGNF\_LIDAR-HD\_MNT\_ELEVATION.ELEVATIONGRIDCOVERAGE.WGS84G} du même
WMS ; le MNS contient les toits et les cimes, on se poserait dessus), en
binaire brut (\code{demanderAltitudes()}). Les trous et les points aberrants
sont bouchés avec le relief de la carte (\code{ReliefCarte::altitude()}), puis
tout est écrit en tuiles \code{.r16} sur 16 bits (\code{ecrireBlocRelief()},
\code{encoderTuile()}). Le pas vient de \code{grilleRelief()} : la maille de
la carte divisée par un multiple de 4, par axe. Un pas libre ferait voir la
frontière de la fenêtre en vol. Le script
\code{tools/terrain/fetch\_relief.py} fait pareil en ligne de commande (le
format des tuiles est décrit dans son en-tête et dans
\code{tools/terrain/relief\_tuiles.py}).

Au chargement, \code{cheminJeuDeRelief()} trouve le dossier
\code{<carte>.relief} et \code{FenetreRelief::ouvrir()}
(\code{src/render/relief/}) alloue une texture d'altitudes avec ses niveaux de
réduction, et deux grilles sans données : le noyau (pas fin, environ 2~m) et
son anneau (pas quadruple). \code{FenetreRelief::suivre()} recentre le tout
sur l'appareil à chaque image.

Au dessin (\code{Application::renderTerrainAndBuildings()}), l'anneau est tiré
d'abord, le noyau par-dessus (profondeur \code{GL\_LEQUAL} : à égalité, le
dernier tiré gagne). Le pochoir écarte le maillage d'ensemble de l'emprise de
la fenêtre. Dans \code{terrain.vert}, chaque sommet retrouve sa position
depuis son numéro (\code{gl\_VertexID}), lit l'altitude de la carte
(\code{hauteurCarte()}) et l'écart du laser (\code{detailFin()}, au niveau de
réduction donné par \code{finesseDetail()}), puis fond cet écart vers zéro au
bord de la fenêtre (\code{poidsFenetre()}) : \code{hauteurFenetre()} = carte +
poids $\times$ détail. À la frontière, la surface rejoint donc exactement le
maillage d'ensemble. Les normales sont recalculées par différences finies
(\code{normaleDe()}), et \code{terrain.frag} écarte les points hors de
l'emprise de la carte.

La physique suit la même formule : \code{FenetreRelief::detailEn()} recalcule
côté processeur le même écart et le même poids, \code{Terrain::heightAt()}
l'ajoute à l'altitude de la carte. On se pose sur ce qu'on voit ; c'est un
invariant du moteur.

\section{Modélisation}

En un point $p$ du sol, à distance $d$ de l'oeil de la fenêtre :
\begin{equation*}
  h(p) = c(p) + t(d)\,\bigl[\,T_{\ell}(p) - c(p)\,\bigr]
\end{equation*}
$c(p)$ est l'altitude de la carte (bilinéaire sur \code{heightmap.png}),
$T_{\ell}(p)$ l'altitude laser lue au niveau de réduction $\ell$, $t(d)$ le
poids du fondu. Le détail est l'écart $T - c$, jamais une altitude : quand $t$
tombe à zéro, la surface \emph{est} la carte. Aucune marche possible à la
frontière.

La finesse ne dépend que de la distance :
\begin{equation*}
  \ell(d) = \min\Bigl(\max\bigl(\log_2 \tfrac{d}{D},\, 0\bigr),\, L\Bigr)
  \qquad\text{avec}\qquad
  L = \Bigl\lfloor \log_2 \tfrac{\Delta_{\text{carte}}}{\Delta_{\text{relief}}} \Bigr\rceil
\end{equation*}
Pleine finesse jusqu'à $D$ (environ 226~m sur dax), puis résolution divisée
par deux à chaque doublement de distance, plafonnée à $L$, le niveau qui lisse
le laser à la maille de la carte. Comme $\ell$ ne dépend que de $d$, le noyau
et l'anneau lisent la même surface là où ils se rejoignent.

Le pas des tuiles s'emboîte dans la maille $m$ de la carte :
\begin{equation*}
  \Delta_{\text{relief}} = \frac{m}{4k},
  \qquad k \in \mathbb{N}^{*} \text{ choisi pour approcher } 2~\text{m}
\end{equation*}
Sur dax : $m = 14{,}156 \times 21{,}763$~m donne $14{,}156/8 = 1{,}7695$~m et
$21{,}763/12 = 1{,}8136$~m. 8 et 12 sont multiples de 4 : le noyau et l'anneau
(qui dessine à 4 fois ce pas) tombent tous deux sur les noeuds du maillage.

Chaque tuile de $512 \times 512$ points quantifie ses altitudes sur 16 bits,
bornes propres à la tuile :
\begin{equation*}
  a = a_{\min} + \frac{n}{65535}\, e,
  \qquad n \in \{0, \dots, 65535\},
  \qquad \varepsilon_{\max} = \frac{e}{131070}
\end{equation*}
où $e$ est l'étendue d'altitude de la tuile : 1~cm d'erreur sur 1300~m de
dénivelé, 0,08~mm sur une tuile de plaine de 10~m d'étendue. Mesuré contre la
chaîne Python de référence : 0,59~mm d'écart maximal sur 2,4 millions de
points.

Les normales sont des différences finies centrales, au pas $\delta$ suivant la
finesse :
\begin{equation*}
  \vec{n} = \operatorname{normalize}\!\left(
  \frac{h(x-\delta,z) - h(x+\delta,z)}{2\delta},\;
  1,\;
  \frac{h(x,z-\delta) - h(x,z+\delta)}{2\delta}
  \right)
\end{equation*}

Pourquoi la double précision : autour de 43 degrés, deux \code{float}
consécutifs sont séparés de
\begin{equation*}
  \operatorname{ulp}(43) = 43 \times 2^{-23}
  \approx 5{,}1 \times 10^{-6}~\text{degré}
  \approx 40~\text{cm au sol.}
\end{equation*}
En simple précision, le relief sortait décalé d'un demi-mètre sur les
versants. En double, le même écart tombe au dixième de micron.

\section{Voir la mécanique en vol}

La clé \code{relief\_debug 1} de \code{assets/config.txt} trace les frontières
de la fenêtre : contour du noyau en magenta, contour de l'anneau en cyan,
cercles du fondu en orange et jaune, légende dans le HUD 4 coins (touche
\code{H}). Le journal du lancement donne le reste : maille LR, finesse des
tuiles HR, pas des grilles, et la ligne \code{emboîtement dans la carte : oui}.

\section{Les difficultés rencontrées, et les solutions}

\textbf{Des silhouettes qui bougeaient à la frontière de la fenêtre.} Avec un
pas de tuiles de 2,00~m partout, la fenêtre redessinait la surface du maillage
au lieu de la reproduire : les crêtes changeaient de forme en entrant.
Solution : le pas n'est plus libre, \code{grilleRelief()} le cale sur la
maille de la carte divisée par un multiple de 4, par axe
(\code{pasEmboite()}).

\textbf{Des pics qui surgissaient à l'approche des crêtes.} En réalité un
peigne de lames verticales, une par cellule d'anneau, à partir du cercle de
départ du fondu. Les données, la loi de finesse, la couture de la texture et
les indices des grilles ont été mis hors de cause un par un, mesure après
mesure. La cause : le noyau, tiré en premier, réservait ses pixels au
pochoir ; sur une arête vue en rasant, ses triangles de 2~m se projettent en
lamelles de moins d'un pixel, les pixels non couverts restaient libres et
l'anneau y montrait sa propre corde, plus basse. Chaque grille seule dessinait
la crête proprement ; le défaut venait de leur arbitrage. Solution : inverser
l'ordre, l'anneau d'abord sur toute son emprise, le noyau par-dessus en
\code{GL\_LEQUAL} (à égalité, le dernier tiré gagne). Coût nul. La physique
n'a jamais été touchée : \code{detailEn()} lit le champ continu, c'est pour ça
que l'appareil s'enfonçait dans les dents dessinées.

\textbf{Des changements de couleur en approchant une montagne.} Deux causes.
Au bord de carte, la fenêtre débordait de l'emprise et ses textures en
répétition y dessinaient l'orthophoto et le relief de l'autre côté de la
carte : corrigé par un écart au fragment hors emprise (\code{terrain.frag}).
Et un versant pouvait passer clair, sombre, puis clair pendant l'approche :
les normales sont recalculées à un pas qui suit la finesse
(\code{pasNormale()}), l'éclairage dépend donc de la distance. L'essentiel a
disparu avec le peigne ; le reste est surveillé.

\textbf{Un décalage de 80~cm du relief sur les versants.} Les bornes
géographiques de la fabrique étaient en simple précision (voir la section des
formules). Passées en double (\code{CalageCarte}), ce qui a aussi corrigé le
même décalage sur les tuiles d'orthophoto.

\textbf{Des points laser qui ne mesuraient pas le sol.} Certains retours
plongent de plusieurs centaines de mètres sous le relief (jusqu'à $-796$~m sur
un versant d'Ossau). La fabrique écarte tout point à plus de 300~m sous la
carte (\code{CHUTE\_ABERRANTE\_M}) et bouche trous et aberrations avec le
relief d'ensemble (\code{ReliefCarte::altitude()}).

À retenir de l'enquête du peigne : n'afficher qu'une grille à la fois pour
isoler le fautif, compter les traversées de silhouette le long d'une ligne
plutôt que des écarts entre pixels, couper le MSAA avant toute sonde qui
encode des mesures dans les couleurs.

\newpage

\section{Sources de données}

\begin{itemize}
  \item Géoplateforme IGN, service WMS raster :
    \url{https://data.geopf.fr/wms-r/wms} (orthophotos BD ORTHO et MNT LiDAR
    HD).
  \item API altimétrie IGN (RGE ALTI) :
    \url{https://data.geopf.fr/altimetrie/1.0/calcul/alti/rest/elevation.json}.
  \item Portail des géoservices : \url{https://geoservices.ign.fr/}.
  \item Licence Ouverte Etalab 2.0 :
    \url{https://www.etalab.gouv.fr/licence-ouverte-open-licence/}.
\end{itemize}

\end{document}
